% !TeX root = ../nc-del-pezzo.tex
% 2026-11-20 JST / v0.11
\section{全方向有限族の低次数中心と通常次数付き表示の認識}
\label{sec:simultaneous-low-center}

固定次数の中心条件だけを用いることで，形式一般例にとどまらず，
全方向有限族の非空稠密な開集合に
定理\ref{thm:ncdp-veronese-lift}を適用できる。
さらに方向解消のproper性により，この低次数中心条件は
可換零切断全体を含む一つの開近傍の非零部分で確保できる。
本節では全次数の中心が最小となる一つのZariski開集合の存在を主張しない。

定理\ref{thm:simultaneous-degree-one-central-family}または
定理\ref{thm:simultaneous-degree-two-central-family}の有限基底を $X_d$ と書く。
中心元の次数を $e=6$ または $e=4$ とし，
中心元の構成に用いた完全対称化
\[
 (\operatorname{Sym}V)_e\otimes\mathcal O_{X_d}\longrightarrow A_e
\]
が同型となる開集合への制限を含める。
ここで $V$ はそれぞれ重み $(1^n,2,3)$ または $(1^{n+1},2)$ の生成空間である。
さらに必要に応じて，定理\ref{thm:simultaneous-smooth-ncdp-open}の
滑らかな可換零切断全体を含む開集合 $V_d\subset X_d$ へ制限してよい。
方向が非零である部分を $X_d^\times$，$V_d^\times$ と書く。
その上では
\[
 \mathcal K=
 \begin{cases}
  \ker(D:E\otimes\mathcal O\to\mathcal O),&d=1,\\
  \ker(P:E\otimes\mathcal O\to E^*\otimes\mathcal O),&d=2
 \end{cases}
\]
は階数 $k=n-1$ のベクトル束である。
表示関係式からその切断は厳密に中央であり，
$\mathcal P=\operatorname{Sym}_{\mathcal O}\mathcal K$ と置く。

\subsection{中央多項式環上の普遍的平坦性}

\begin{proposition}\label{lem:simultaneous-central-flatness}
$X_d^\times$ 上の同じ中心商 $B=A/(\Omega)$ について，
$\mathcal P\to B$ は単射であり，$B$ は $\mathcal P$ 上平坦である。
より正確には
\[
 C=B/(\mathcal K)B
\]
の各 $C_m$ は有限局所自由であり，基底の各affine開集合上で
非標準的な次数付き $\mathcal P$-加群同型
\begin{equation}\label{eq:simultaneous-casimir-splitting}
 B\simeq\mathcal P\otimes_{\mathcal O}C
\end{equation}
を取れる。この同型では $C_0=\mathcal O$ の単位元を $1\in B$ へ送れる。
これらの単射性，平坦性，次数別有限局所自由性は任意の可換基底変更後も成立する。
特に $B_m/\mathcal P_m$ は有限局所自由で，その形成は任意基底変更と可換する。
\end{proposition}
\begin{proof}
\emph{閉ファイバーの活性最高次元。}
まず任意の複素点 $x\in X_d^\times$ を固定し，
$\mathcal K_x$ の基底を $s_1,\ldots,s_k$ とする。
活性変数を次数1では $(u,y,z)$，次数2では $(u,v,z)$ とし，
その活性重みをそれぞれ $(1,2,3)$，$(1,1,2)$ とする。
中央変数 $s_i$ の活性重みは0とするが，内部次数は引き続き1である。
定理\ref{thm:simultaneous-poisson-ambient}および
定理\ref{thm:simultaneous-degree-one-ambient}のPBW証明から
\[
 \operatorname{gr}_{\mathrm{act}}A_x
   \simeq P_x\otimes_{\mathbb C}A_x^{\mathrm{top}},
 \qquad P_x=\mathbb C[s_1,\ldots,s_k]
\]
を得る。$A_x^{\mathrm{top}}$ はそれぞれの全最高次AS開条件を満たす整域であり，
標準の三変数Hilbert級数を持つ。
次数2の整域性には\cite[Corollary~3.8]{MoriUeyama2019}を用いる。
この議論は活性最高次曲線の滑らかさを要求しない。

$\Omega_x$ の活性最高次が文字通り $e$ であることを確認する。
内部次数 $e$ の完全対称化写像は上の活性filtrationを保つ。
その始域と終域について，各filtration部分の次元はPBW同型により等しい。
従って，filtrationに適合した基底で完全対称化を表す行列は，
正方な対角ブロックを持つ三角行列である。
全体の写像は $X_d$ の選択により同型なので，全対角ブロックも同型である。
特に活性次数 $e$ のブロック
\[
 \mathbb C[\text{活性変数}]_e
      \longrightarrow (A_x^{\mathrm{top}})_e
\]
は同型である。
完全対称化座標で $\Omega_x$ の $z^2$ 係数は $1/2$ であり，
$z^2$ の活性次数は $e$ だから，その活性最高記号
$\omega_x\in(A_x^{\mathrm{top}})_e$ は零でない。
内部斉次性から，この最高記号に $s_i$ は現れない。

$\Omega_x$ の中心性より $\omega_x$ も中央である。
$A_x^{\mathrm{top}}$ は整域なので，$\omega_x$ は左右正則である。
最高記号による次数降下を中央元の掛算に適用すると，
その掛算はstrictであり，
\begin{equation}\label{eq:simultaneous-casimir-fiber-gr}
 \operatorname{gr}_{\mathrm{act}}B_x
  \simeq P_x\otimes_{\mathbb C}Q_x,
 \qquad Q_x=A_x^{\mathrm{top}}/(\omega_x)
\end{equation}
となる。
各 $(Q_x)_j$ の基底を内部次数 $j$ の元へ持ち上げれば，
活性次数の帰納法により，それらは $B_x$ の斉次 $P_x$-自由基底となる。
従って $B_x$ は $P_x$ 上自由であり，
$C_x=B_x/(s_1,\ldots,s_k)B_x$ のHilbert級数は
\begin{equation}\label{eq:simultaneous-active-quotient-hilbert}
 H_{C_x}(t)=
 \begin{cases}
 \dfrac{1-t^6}{(1-t)(1-t^2)(1-t^3)},&d=1,\\[6pt]
 \dfrac{1-t^4}{(1-t)^2(1-t^2)},&d=2.
 \end{cases}
\end{equation}
特にこの級数は $x$ に依存しない。

\emph{普遍係数上の有限局所自由性。}
基底のaffine開集合 $\operatorname{Spec}S\subset X_d^\times$ で作業する。
$S$ は被約な有限型 $\mathbb C$-代数である。
各 $C_m$ は $B_m$ の有限表示商であり，商の形成は基底変更と可換する。
上で全閉ファイバーの次元が一定であることを示したので，
有限表示行列の一定階数判定により $C_m$ は有限局所自由となる。
ここで被約性を使う箇所を明記する。
関係行列の所要階数より大きい小行列式は全閉点で零となるから，
有限型被約環 $S$ 上で恒等的に零である。
所要階数の非零小行列式を局所化して消去すれば残る行列は零となり，
$C_m$ の自由性を得る。
非被約基底上でファイバー次元だけからこの結論を推論してはいない。

各 $C_m$ は $S$ 上有限射影なので，$B_m\twoheadrightarrow C_m$ の
$S$-線型な分裂を次数ごとに選べる。
次数0では $1$ を $1$ へ持ち上げる。
これらを掛算により延長すると，次数付き $P=\operatorname{Sym}_S\mathcal K$-線型写像
\[
 \mu:P\otimes_SC\longrightarrow B
\]
を得る。
内部次数の帰納法，すなわち正次数ideal $(\mathcal K)$ に関する
次数付きNakayamaの議論により，$\mu$ は全射である。
各次数の始域と終域は有限射影 $S$-加群であり，
全閉ファイバーの階数は式\eqref{eq:simultaneous-active-quotient-hilbert}と
$H_{B_x}=(1-t)^{-k}H_{C_x}$ により等しい。
従って各 $\mu_m$ は同型であり，式\eqref{eq:simultaneous-casimir-splitting}を得る。
この同型は $P$ を $C_0=S$ に対応する直和因子へ送るため，$P\to B$ も単射である。

最後に式\eqref{eq:simultaneous-casimir-splitting}そのものを任意の $S$-代数へ
基底変更すれば，非被約な基底を含めて主張が保たれる。
$C=\bigoplus_{m\ge0}C_m$ は射影加群の直和なので，
$P\otimes_SC$ は $P$ 上平坦である。
各固定次数には有限個の直和因子だけが寄与することから，
$B_m/P_m$ の有限局所自由性と任意基底変更との両立も従う。
\end{proof}

\subsection{一つの固定次数における普遍中心の稠密開集合}

\begin{theorem}\label{thm:finite-low-center-open}
正整数 $q$ を一つ固定する。
$X_d^\times$ には有限行列の最大小行列式で定義される座標変換に不変な開集合 $U_{d,q}$ があり，
任意の基底変更 $T\to U_{d,q}$ に対し
\begin{equation}\label{eq:simultaneous-universal-low-center}
 Z(B_T)_q=(\operatorname{Sym}_{\mathcal O_T}\mathcal K_T)_q
\end{equation}
が成り立つ。
滑らかな可換方程式の係数 $f$ と任意の非零Poisson方向 $\xi$ を固定する。
有限族の定義域に含まれる直線
\[
 h\longmapsto(f,h\xi)
\]
の非零部分は $U_{d,q}$ と非空Zariski開集合で交わる。
同じ主張は $V_d^\times$ 内でも成立する。
従って $U_{d,q}\cap V_d^\times$ は非空稠密開集合であり，
有限の非零複素パラメータを持つ例を含む。
\end{theorem}
\begin{proof}
生成元の重みを $w$，その生成空間の自明束を $\mathcal V_w$ とする。
命題\ref{lem:simultaneous-central-flatness}により，有限局所自由加群の写像
\begin{equation}\label{eq:simultaneous-low-center-matrix}
 \delta_q:B_q/\mathcal P_q\longrightarrow
       \bigoplus_w\operatorname{Hom}_{\mathcal O}
          (\mathcal V_w,B_{q+w}),
 \qquad [b]\longmapsto(v\mapsto[b,v])
\end{equation}
が定義される。
これらは代数生成元全てとの交換子なので，中心の条件はその核が零であることに等しい。
幾何学的ファイバーで $\delta_q$ が単射となるlocusを $U_{d,q}$ とする。
これは有限な最大小行列式による開条件である。
$\mathcal K$ と交換子写像の構成は座標変換に自然なので，この開集合は不変である。
この開集合上では $\delta_q$ は局所分裂単射となるので，
任意の非被約基底変更後も単射であり，
式\eqref{eq:simultaneous-universal-low-center}を得る。

非零方向ごとの非空性には，直後の定理\ref{thm:low-center-zero-neighborhood}を用いる。
その証明は，ここで定義した交換子行列と普遍平坦性だけに依存し，
本定理の非空性の結論を使わない。
同定理が与える零切断を含む開近傍 $W_{d,q}$ の
直線 $(f,h\xi)$ への逆像は $h=0$ を含むZariski開集合である。
従ってその非零部分は非空であり，その全体が $U_{d,q}$ に含まれる。
$V_d$ の逆像との交叉についても同じである。
最後に $V_d^\times$ はaffine空間，またはaffine空間と
既約Pl\"ucker錐の積の非空開であるから既約であり，
得られた非空開集合は稠密である。
\end{proof}

\subsection{方向解消による零切断近傍での低次数中心条件}

以下では，全方向有限族の基底 $V_d$ を，必要なら可換方程式の滑らかな
係数開集合 $\mathcal T_d$ の逆像へ制限する。
この操作は滑らかな可換模型の零切断全体
\[
 Z_d=\mathcal T_d\times\{0\}
\]
を保つ。
従って $V_d$ は $\mathcal T_d\times\mathcal C_d$ の開部分であり，
$\mathcal C_1=E^*$，$\dim E=n$，
$\mathcal C_2\subset\bigwedge^2E^*$ はrankが高々2のPl\"ucker錐，
$\dim E=n+1$ である。
$B$ は既存の同時中心商で，各 $B_m$ は有限局所自由である。
$V_d^\times=V_d\setminus Z_d$ と置く。

\begin{theorem}[零切断全体を保つ低次数中心の一様化]
\label{thm:low-center-zero-neighborhood}
正整数 $q$ を固定する。
$Z_d$ 全体を含むZariski開集合 $W_{d,q}\subset V_d$ が存在し，
\[
 W_{d,q}\setminus Z_d\subset U_{d,q}
\]
となる。ここで $U_{d,q}$ は
定理\ref{thm:finite-low-center-open}の有限行列で定義した開集合である。
特に任意の可換基底変更 $T\to W_{d,q}\setminus Z_d$ に対して
\begin{equation}\label{eq:zero-neighborhood-low-center}
 Z(B_T)_q=\operatorname{Sym}^q_{\mathcal O_T}\mathcal K_T
\end{equation}
が成り立ち，$B_T$ は
$\operatorname{Sym}_{\mathcal O_T}\mathcal K_T$ 上平坦である。
$V_d$ を座標変換に不変に選べば，$W_{d,q}$ も不変に選べる。

この結論は一つの固定次数 $q$ についてのものである。
零切断そのものでは $B$ は可換であり，
式\eqref{eq:zero-neighborhood-low-center}を $W_{d,q}$ 全体へ延長してはいない。
\end{theorem}
\begin{proof}
\emph{全係数を保ったproperな方向解消。}
方向の射影化とそのtautological lineを
\[
 \Gamma_d=
 \begin{cases}
  \mathbb P(E^*),&d=1,\\
  \operatorname{Gr}(2,E^*),&d=2,
 \end{cases}
 \qquad \mathcal L=\mathcal O_{\Gamma_d}(-1)
\]
とする。
次数1では標準評価列
\[
 0\longrightarrow\mathcal K\longrightarrow
 E\otimes\mathcal O_{\Gamma_1}\longrightarrow\mathcal L^\vee
 \longrightarrow0
\]
を，次数2では標準商束 $Q$ の列
\[
 0\longrightarrow\mathcal K\longrightarrow
 E\otimes\mathcal O_{\Gamma_2}\longrightarrow Q\longrightarrow0,
 \qquad \operatorname{rank}Q=2,
 \quad\mathcal L=\det Q^*
\]
を使う。いずれも $\mathcal K$ の階数は $n-1$ である。
閉incidence表示から，方向の射
\[
 \operatorname{Tot}(\mathcal L)\longrightarrow\mathcal C_d
\]
はproperで，頂点外では同型である。
これを係数空間と積にし，$V_d$ へ基底変更して
\[
 \nu:\widetilde V_d=
  (\mathcal T_d\times\operatorname{Tot}(\mathcal L))
       \times_{\mathcal T_d\times\mathcal C_d}V_d
       \longrightarrow V_d
\]
を得る。
例外因子は
\[
 D_d=\nu^{-1}(Z_d)=\mathcal T_d\times\Gamma_d
\]
である。$Z_d$ 全体が $V_d$ に入るため，例外因子から方向や係数が失われない。

以下，$\mathcal L,\mathcal K$ の引戻しも同じ記号で書く。
線束全空間のtautological sectionを $\tau$ とすると，
$D_d=(\tau=0)$ は有効Cartier因子で，
$\mathcal I_{D_d}\simeq\mathcal L^\vee$ である。
局所フレーム $\ell$ を選んで $\tau=h\ell$ と書けば，
普遍方向は $D=hD_0$ または $P=hP_0$ と書ける。
ここで $D_0$ は非零線型形式，$P_0$ はrank2の交代形式であり，
それぞれ $\mathcal K$ を核に持つ。
$\widetilde B=\nu^*B$ の各次数は有限局所自由なので，
$h$ の掛算は各次数で単射である。
方向解消の射 $\nu$ 自体の平坦性は用いない。

\emph{中央部分の商が全解消上でベクトル束になること。}
関係式\eqref{eq:simultaneous-degree-one-relations}と
\eqref{eq:simultaneous-ambient-relations}から，
$\mathcal K$ の切断は $h=0$ も含めて厳密に中央である。
従って積による写像
\[
 \iota_q:\operatorname{Sym}^q\mathcal K\longrightarrow\widetilde B_q
\]
がある。
例外因子の点 $(f,[\xi])$ では，$\mathcal K$ に適合する座標を選ぶと
可換切断環は $P=\mathbb C[s_1,\ldots,s_{n-1}]$ 上，
活性二変数の単項式と $z$ の指数0または1を基底とする自由加群である。
これは方程式が $z$ に関してmonicであることによる。
この基底は $1$ を含むので，$\iota_q$ のファイバーは単射である。
例外因子外では，命題\ref{lem:simultaneous-central-flatness}の
普遍的な中央多項式加群の分裂により同じ単射性が成立する。
従って $\iota_q$ は $\widetilde V_d$ の全ファイバーで単射である。

有限局所自由加群の間の写像が全ファイバーで単射なら，
適切な最大サイズの小行列式を可逆化して局所的な左逆を取れる。
従って $\iota_q$ は局所分裂単射であり，
\begin{equation}\label{eq:resolved-low-degree-bundle}
 0\longrightarrow\operatorname{Sym}^q\mathcal K
  \xrightarrow{\iota_q}\widetilde B_q
  \longrightarrow\mathcal Q_q\longrightarrow0
\end{equation}
の商 $\mathcal Q_q$ は $\widetilde V_d$ 全体で有限局所自由である。
この完全列と商の形成は任意の基底変更に両立する。
この箇所ではファイバー次元だけから非被約基底上の平坦性を推論せず，
有限束の写像の最大階数による局所分裂を用いている。

\emph{除した交換子と任意基底変更。}
重み $w$ の固定された生成空間を $V_w$ とし，
\[
 \mathcal M_q=
 \bigoplus_w\operatorname{Hom}_{\mathcal O_{\widetilde V_d}}
 (V_w\otimes\mathcal O_{\widetilde V_d},\widetilde B_{q+w})
\]
と置く。
生成元全てとの交換子の写像
$\delta_q:\widetilde B_q\to\mathcal M_q$ は，
$D_d$ 上の代数が可換なので $D_d$ への制限で零となる。
従って
\[
 0\longrightarrow\mathcal M_q\otimes\mathcal L^\vee
   \xrightarrow{\ \tau\ }\mathcal M_q
   \longrightarrow\mathcal M_q|_{D_d}\longrightarrow0
\]
を用いて，一意的な写像
\[
 \mathcal D_q:\widetilde B_q\longrightarrow
                    \mathcal M_q\otimes\mathcal L^\vee,
 \qquad \tau\mathcal D_q=\delta_q
\]
を得る。局所的には $\mathcal D_q=h^{-1}\delta_q$ である。
フレームを $\ell'=a\ell$ に替えると $h'=a^{-1}h$ なので，
局所式は $h'^{-1}\delta_q=a h^{-1}\delta_q$ と変化し，
表示した $\mathcal L^\vee$ の捩れと一致する。
$\operatorname{Sym}^q\mathcal K$ は厳密に中央であり，
$\tau$ の掛算は単射なので，$\mathcal D_q$ はこれを零に送る。
従って商束上の有限行列
\begin{equation}\label{eq:resolved-divided-commutator}
 \overline{\mathcal D}_q:\mathcal Q_q
    \longrightarrow\mathcal M_q\otimes\mathcal L^\vee
\end{equation}
が得られる。

\emph{例外因子の全点における核。}
任意の $(f,[\xi])\in\mathcal T_d\times\Gamma_d$ を取る。
関係式から，$\mathcal D_q$ のこの点でのファイバーは
方向 $\xi$ の古典Poisson括弧による写像
\[
 (B_f)_q\longrightarrow
       \bigoplus_w\operatorname{Hom}(V_w,(B_f)_{q+w}),
 \qquad b\longmapsto\bigl(v\longmapsto\{b,v\}_\xi\bigr)
\]
である。ただし $\xi$ の非零スカラー倍による変更は
値域の一次元因子を変えるだけなので，核に影響しない。
これは同じ有限表示族の一次交換子の計算であり，
別の量子化との形式的同型からの推論ではない。
また $\Omega|_{D_d}=F_f$ なので，括弧はまさに可換切断環 $B_f$ 上の括弧である。

補題\ref{lem:classical-casimir-center}は，
\emph{任意の滑らかな係数 $f$ と任意の非零方向 $\xi$} に対して
\[
 Z_{\mathrm{Pois}}(B_f)=\operatorname{Sym}\mathcal K_{[\xi]}
\]
を与える。
その適用範囲をここで再確認する。
$P=\mathbb C[s_1,\ldots,s_{n-1}]$，$K=\operatorname{Frac}P$ とする。
次数1では，活性変数が $(u,y,z)$ であり，$z$ の二次式の定数項は
$y$ について次数3を持つので，$\overline K[u,y]$ の平方にはならない。
従って $B_f\otimes_PK$ は幾何学的整である。
次数2で活性四次式が $\overline K[u,v]$ の平方なら，
平方根の非零係数を1に正規化してGalois降下することにより
$f_4=a g^2$，$a\in K^\times$，$g\in K[u,v]$ と書ける。
$g$ が活性変数に関して非定数なら，Gaussの補題から
$P[u,v]$ における活性変数を含む因子が重複し，
分岐四次超曲面の滑らかさに反する。
$g$ が定数なら $f_4\in P$ であり，
$s_1=\cdots=s_{n-1}=0$ で定まる活性射影直線上で
$f_4$ と全偏微分が消えるので，やはり滑らかでない。
よって次数2でも同じ幾何学的整性を得る。

両次数で活性な最初の二変数の括弧は $z$ の非零定数倍であり，
$\operatorname{Frac}(B_f)/K$ は超越次数2のPoisson体でgeneric rank2を持つ。
従って体のPoisson中心は $K$ 上代数的な元からなり，
上の幾何学的整性からちょうど $K$ に等しい。
monicな $z$ 除法の $P$-自由基底が $1$ を含むため，
$B_f\cap K=P$ であり，環のPoisson中心は $P$ となる。
この議論には活性最高次曲線の滑らかさも，係数 $f$ の一般性も不要である。
代数閉拡大体上でも同じ証明が成立する。

したがって式\eqref{eq:resolved-divided-commutator}は
例外因子の全幾何学的点で単射である。
この写像がファイバー上単射となる開集合を
$\widetilde W_{d,q}\subset\widetilde V_d$ とすれば，
\[
 D_d\subset\widetilde W_{d,q}
\]
である。
この開集合上では $\overline{\mathcal D}_q$ は局所分裂単射なので，
式\eqref{eq:resolved-low-degree-bundle}と合わせて
\[
 \ker(\mathcal D_{q,T})=\operatorname{Sym}^q\mathcal K_T
\]
が任意の基底変更 $T\to\widetilde W_{d,q}$ の後も成立する。
ここで $\mathcal D_{q,T}$ は，元の基底で定義した除した交換子の引戻しである。
$\tau_T$ が零因子となる基底上で新たに割算をしているのではない。

\emph{proper性による対象側の開近傍。}
悪い集合 $C_{d,q}=\widetilde V_d\setminus\widetilde W_{d,q}$ は閉であり，
$\nu$ がproperなので $\nu(C_{d,q})$ も $V_d$ の閉集合である。
さらに $C_{d,q}\cap\nu^{-1}(Z_d)=\varnothing$ なので，
\[
 W_{d,q}=V_d\setminus\nu(C_{d,q})
\]
は $Z_d$ 全体を含む開集合である。
この構成は全係数空間 $\mathcal T_d$ の上で行っており，
固定した一つの $f$ ごとに得た開集合の交叉を取ってはいない。

$\nu^{-1}(W_{d,q})\subset\widetilde W_{d,q}$ であり，
頂点外では $\nu$ は同型，$\tau$ は可逆である。
従って $W_{d,q}\setminus Z_d$ 上では通常の交換子写像
$\delta_q$ の商束上の写像も局所分裂単射となる。
これは $W_{d,q}\setminus Z_d\subset U_{d,q}$ を意味し，
任意基底変更後の中心の式を与える。
中央多項式環上の平坦性は命題\ref{lem:simultaneous-central-flatness}による。
全構成は座標変換に自然なので，不変性も従う。
\end{proof}


\begin{corollary}\label{cor:finite-full-graded-recognition}
$n\ge3$，$r=n-1$ とし，$U=U_{d,r}\cap V_d^\times$ と置く。
任意の基底 $T$ 上の二つの点 $a,b:T\to U$ とその中心商 $B_a,B_b$ に対して，
制限写像
\[
 \underline{\operatorname{Isom}}_{\mathrm{gr},T}(B_a,B_b)
 \longrightarrow
 \underline{\operatorname{Isom}}_{\mathrm{gr},T}
           (B_a^{(r)},B_b^{(r)})
\]
はfppf位相で $\mu_{r,T}$-torsorであり，
\[
 \underline{\operatorname{Isom}}_{\mathrm{gr},T}(B_a,B_b)/\mathbb G_m
  \simeq
 \underline{\operatorname{Isom}}_{\mathrm{gr},T}
           (B_a^{(r)},B_b^{(r)})/\mathbb G_m
\]
である。両辺のスカラー作用の対応は $\lambda\mapsto\lambda^r$ である。
さらに定理\ref{thm:finite-anticanonical-isom}により，これらの層は
\begin{equation}\label{eq:finite-full-graded-recognition}
 \underline{\operatorname{Eq}}_{\mathcal O}(\QGr B_a,\QGr B_b)
 \simeq\underline{\operatorname{Isom}}_{\mathrm{gr},T}(B_a,B_b)/\Gm
\end{equation}
と同一視できる。これは $U\times U$ 上affineかつ有限表示である。
$V_d$ を系\ref{cor:finite-equivalence-unramified}の開集合に取れば，
この同値スキームはquasi-finiteかつunramifiedでもある。
\end{corollary}
\begin{proof}
命題\ref{lem:simultaneous-central-flatness}が中央多項式環上の平坦性を，
定理\ref{thm:finite-low-center-open}が任意基底変更後の次数 $r$ の中心条件を与える。
$\operatorname{rank}\mathcal K=r\ge2$ なので，
定理\ref{thm:ncdp-veronese-lift}の全代数的仮定を満たす。
圏論的結論と表現可能性には定理\ref{thm:finite-anticanonical-isom}を用いる。
\end{proof}

\begin{remark}
$n=2$ では $r=1$ なので，Veronese持上げは不要であり，
通常の $B$ による認識は定理\ref{thm:finite-anticanonical-isom}の $V_d$ 全体で成立する。
$n\ge3$ で使った中心条件は $q=r$ という一つの有限次数だけであり，
全中心の最小性や全モジュライのatlasを結論していない。
また通常の次数付き代数同型では，重み2以上の生成元に低い重みの生成元の
同次数多項式を加える変換も許す。
そのような変換を除外した線型生成元変換群だけで同型層を代用してはいけない。
\end{remark}
